\documentclass[12pt,a4paper]{article}
\usepackage[utf8]{inputenc}
\usepackage[T1]{fontenc}
\usepackage[french]{babel}
\usepackage{geometry}
\usepackage{titlesec}
\usepackage{booktabs}
\usepackage{array}
\usepackage{tabularx}
\usepackage{enumitem}
\usepackage{graphicx}
\usepackage{hyperref}
\usepackage{fancyhdr}
\usepackage{xcolor}
\usepackage{parskip}

\geometry{top=2.5cm, bottom=3cm, left=2.5cm, right=2.5cm}

\definecolor{primary}{RGB}{0,70,127}
\hypersetup{
  colorlinks=true,
  linkcolor=primary,
  urlcolor=primary
}

\pagestyle{fancy}
\fancyhf{}
\renewcommand{\headrulewidth}{0.4pt}
\renewcommand{\footrulewidth}{0.4pt}
\fancyhead[L]{\small UE : ANI-IA 4067 (AIA4)}
\fancyhead[R]{\small ENSPY -- Génie Informatique}
\fancyfoot[L]{\small NEUSSI NJIETCHEU Patrice Eugene -- 24P820}
\fancyfoot[C]{\thepage}
\fancyfoot[R]{\includegraphics[height=0.8cm]{enspy.png}}

\fancypagestyle{plain}{
    \fancyhf{}
    \fancyfoot[L]{\small NEUSSI NJIETCHEU Patrice Eugene -- 24P820}
    \fancyfoot[C]{\thepage}
    \fancyfoot[R]{\includegraphics[height=0.8cm]{enspy.png}}
    \renewcommand{\footrulewidth}{0.4pt}
    \renewcommand{\headrulewidth}{0pt}
}

\titleformat{\section}{\large\bfseries\color{primary}}{}{0em}{\thesection\quad}[\titlerule]
\titleformat{\subsection}{\normalsize\bfseries}{}{0em}{\thesubsection\quad}

\setlist[itemize]{noitemsep, topsep=4pt, leftmargin=1.5em}
\setlist[enumerate]{noitemsep, topsep=4pt, leftmargin=1.8em}

\begin{document}

%---------------------------------------------------------------
%  PAGE DE TITRE
%---------------------------------------------------------------
\begin{titlepage}
  \centering
  \includegraphics[width=4cm]{enspy.png}\\[0.5cm]
  {\Large\bfseries ÉCOLE NATIONALE SUPÉRIEURE POLYTECHNIQUE DE YAOUNDÉ\par}
  \vspace{0.2cm}
  {\normalsize Département de Génie Informatique\par}
  {\normalsize Cycle des Ingénieurs du Numérique -- Filière : AIA4\par}
  \vspace{2cm}
  
  {\LARGE\bfseries\color{primary} PROPOSITION DE PLAN DE TRAVAIL\par}
  \vspace{0.5cm}
  {\large Démarches d'Implémentation et Méthodologies de Mise en Œuvre des Projets d'IA\par}
  \vspace{0.2cm}
  {\large\bfseries UE : ANI-IA 4067 INTELLIGENCE ARTIFICIELLE (AIA4)\par}
  \vspace{2.5cm}
  
  \begin{tabular}{rl}
    \textbf{Rédigé par :} & \textbf{NEUSSI NJIETCHEU Patrice Eugene} \\[4pt]
    \textbf{Matricule :}    & \textbf{24P820} \\[4pt]
    \textbf{Classe :}       & \textbf{AIA4} (Intelligence Artificielle) \\[4pt]
    \textbf{Sous la direction de :} & \textbf{M. Bitha Junior} (Enseignant ENSPY) \\
  \end{tabular}
  \vfill
  {\small Proposition soumise le 23 avril 2026\par}
\end{titlepage}

%---------------------------------------------------------------
%  TABLE DES MATIÈRES
%---------------------------------------------------------------
\tableofcontents
\newpage

%---------------------------------------------------------------
%  INTRODUCTION
%---------------------------------------------------------------
\section{Introduction et Méthodologie Commune}

Ce rapport présente la proposition de plan de travail détaillé et les démarches méthodologiques envisagées pour mener à bien les trois projets pratiques d'Intelligence Artificielle qui nous ont été alloués dans le cadre du cours ANI-IA 4067. 

Afin de garantir des solutions logicielles et algorithmiques performantes et exploitables, nous adopterons une démarche standardisée découpée en six phases structurées, inspirée de la méthodologie industrielle CRISP-DM.

\subsection{Présentation des Projets Alloués}
Les trois projets alloués à réaliser sont :
\begin{enumerate}
    \item \textbf{Projet 1 : Réseaux de neurones convolutifs pour systèmes de recommandations} (Recommandation de films).
    \item \textbf{Projet 2 : Construction d’un moteur de recherche visuel} (Recherche d'images e-commerce par similarité).
    \item \textbf{Projet 3 : Prédiction du comportement des clients avec l'IA dans une banque} (Scoring d'opportunités marketing).
\end{enumerate}

\subsection{Cycle de Réalisation Général}
Pour chaque projet, le pipeline de travail suivra les étapes suivantes :
\begin{itemize}
    \item \textbf{Étape 1 : Analyse des besoins} -- Spécifications des objectifs fonctionnels.
    \item \textbf{Étape 2 : Préparation des données} -- Exploration (EDA), nettoyage et feature engineering.
    \item \textbf{Étape 3 : Modélisation mathématique} -- Choix et implémentation des algorithmes d'apprentissage automatique.
    \item \textbf{Étape 4 : Validation scientifique} -- Calcul des métriques de performance hors-ligne.
    \item \textbf{Étape 5 : Intégration logicielle} -- Développement d'une interface web interactive sous Django.
    \item \textbf{Étape 6 : Déploiement et tests} -- Mise en production sur une plate-forme cloud serverless (Vercel).
\end{itemize}

\newpage

%---------------------------------------------------------------
%  PROJET 1
%---------------------------------------------------------------
\section{Projet 1 -- Système de Recommandations Graph Neural Networks (LightGCN)}

\subsection{Méthode et Solution Proposées}
Pour pallier les insuffisances du filtrage collaboratif classique par factorisation de matrices SVD (qui ne capture que les relations directes), nous prévoyons de modéliser le problème sous forme de graphe relationnel.
\begin{itemize}
    \item \textbf{Modélisation en Graphe Bipartite :} Les interactions d'évaluation (MovieLens) seront représentées par un graphe bipartite non orienté $G = (U \cup I, E)$ où les arêtes relient les utilisateurs aux films.
    \item \textbf{Architecture LightGCN :} Nous utiliserons un réseau neuronal sur graphes (GNN) simplifié, LightGCN, qui supprime les activations non linéaires et les projections de caractéristiques pour ne conserver que la propagation de voisinage par matrice d'adjacence normalisée symétrique :
    \[
    \tilde{\mathbf{A}} = \mathbf{D}^{-1/2} \mathbf{A} \mathbf{D}^{-1/2}
    \]
    \item \textbf{Algorithme d'Apprentissage :} Le modèle sera optimisé avec la perte de classement pairwise BPR (Bayesian Personalized Ranking) afin de classer les films pertinents au-dessus des films non observés.
    \item \textbf{Inférence Temps Réel :} Les plongements latents optimisés seront sérialisés. Nous prévoyons d'implémenter un moteur d'inférence en production léger utilisant NumPy pour calculer les scores en temps réel par simple produit scalaire.
\end{itemize}

\subsection{Plan de Travail Chronologique}
\begin{itemize}
    \item \textbf{Jalons 1 à 3 (Exploration) :} Analyse de la sparcité de la matrice d'interactions MovieLens, nettoyage et construction de l'adjacence globale creuse.
    \item \textbf{Jalons 4 à 6 (Entraînement) :} Écriture sous PyTorch de la couche de convolution de graphe, implémentation du calcul matriciel creux et de la perte BPR.
    \item \textbf{Jalons 7 à 10 (Déploiement) :} Intégration Django, mise en place d'une interface en Vis.js pour visualiser les liens d'interactions et déploiement final sur Vercel.
\end{itemize}

\newpage

%---------------------------------------------------------------
%  PROJET 2
%---------------------------------------------------------------
\section{Projet 2 -- Moteur de Recherche Visuel de Vêtements (CBIR)}

\subsection{Méthode et Solution Proposées}
Nous prévoyons de concevoir une architecture d'extraction d'images hybride et adaptative :
\begin{itemize}
    \item \textbf{Descripteurs Sémantiques (Deep Learning) :} Nous exploiterons le réseau pré-entraîné MobileNetV2 pour sa légèreté. L'extraction de caractéristiques profondes produira des embeddings de 1280 dimensions après Global Average Pooling.
    \item \textbf{Descripteurs de Repli (Colorimétrie/Texture) :} Afin de faire face aux restrictions de taille de package sur Vercel (limitation des fonctions serverless à 15MB compressés interdisant PyTorch), nous développerons un extracteur mathématique purement basé sur NumPy et Pillow. Ce dernier calculera un histogramme joint 3D dans l'espace HSV (128 dimensions) concaténé à une grille de nuances de gris normalisée Z-score (256 dimensions).
    \item \textbf{Recherche de Similarité :} La comparaison des images s'effectuera en mesurant la similarité cosinus entre la requête et les vecteurs de la base de données.
\end{itemize}

\subsection{Plan de Travail Chronologique}
\begin{itemize}
    \item \textbf{Jalons 1 à 2 (Données) :} Sélection et structuration du catalogue d'images statiques (vêtements) et script de prétraitement.
    \item \textbf{Jalons 3 à 5 (Modélisation) :} Implémentation de la fonction d'extraction MobileNetV2 et codage de l'algorithme colorimétrique HSV / Spatial sous NumPy.
    \item \textbf{Jalons 6 à 10 (Intégration) :} Création de l'interface de téléversement (Drag-and-Drop) en Django, écriture du pipeline de recherche adaptatif gérant les fallbacks et déploiement.
\end{itemize}

\newpage

%---------------------------------------------------------------
%  PROJET 3
%---------------------------------------------------------------
\section{Projet 3 -- Prédiction du Comportement des Clients Bancaires}

\subsection{Méthode et Solution Proposées}
Pour le ciblage prédictif des campagnes de marketing direct de la banque (dataset UCI Bank Marketing), nous concevrons un pipeline de classification d'ensemble :
\begin{itemize}
    \item \textbf{Feature Engineering :} Traitement des valeurs manquantes, encodage One-Hot des variables catégorielles et standardisation robuste des variables financières continues (comme le solde moyen).
    \item \textbf{Classification d'Ensemble (Soft Voting) :} Afin de minimiser la variance et de maximiser la robustesse de prédiction, nous prévoyons de combiner trois classificateurs complémentaires :
    \begin{enumerate}
        \item Une Régression Logistique pénalisée L2.
        \item Une Forêt Aléatoire (100 arbres).
        \item Un réseau de neurones artificiels MLP (Perceptron Multicouche) conçu sous PyTorch.
    \end{enumerate}
    \item \textbf{Gestion du déséquilibre :} Les cibles positives n'étant représentées qu'à hauteur de 11.7\%, nous appliquerons des coefficients de pondération inverses lors du calcul des pertes.
\end{itemize}

\subsection{Plan de Travail Chronologique}
\begin{itemize}
    \item \textbf{Jalons 1 à 3 (Exploration) :} Analyse statistique descriptive univariée et multivariée du dataset bancaire, et sélection des variables significatives.
    \item \textbf{Jalons 4 à 6 (Entraînement) :} Construction du pipeline de transformation scikit-learn, écriture du wrapper PyTorch MLP, et optimisation des hyperparamètres par validation croisée.
    \item \textbf{Jalons 7 à 10 (Déploiement) :} Création d'une interface web Django proposant un simulateur d'opportunité client (formulaire unitaire) ainsi qu'un module de traitement de fichier CSV pour les prédictions en lot (batch), puis déploiement sur le Cloud.
\end{itemize}

\newpage

%---------------------------------------------------------------
%  PLANIFICATION TEMPORELLE
%---------------------------------------------------------------
\section{Planification Temporelle et Jalons}

Le tableau ci-dessous décrit la planification des livrables pour les 10 semaines de travail prévues à compter du 23 avril 2026.

\bigskip

\begin{tabularx}{\linewidth}{cXX}
  \toprule
  \textbf{Période} & \textbf{Tâches Planifiées} & \textbf{Livrable attendu} \\
  \midrule
  Semaines 1 -- 2 & EDA, nettoyage des datasets, feature engineering & Rapports d'analyse descriptive et scripts de préparation. \\
  \addlinespace
  Semaines 3 -- 5 & Entraînement des modèles IA locaux (PyTorch/Sklearn) & Code source des algorithmes et courbes d'apprentissage. \\
  \addlinespace
  Semaines 6 -- 7 & Évaluation quantitative et lissage des métriques & Tableaux comparatifs de performance (Recall, ROC-AUC). \\
  \addlinespace
  Semaines 8 -- 9 & Développement des applications web Django & Code source des contrôleurs Django, templates et visuels. \\
  \addlinespace
  Semaine 10 & Optimisation d'inférence NumPy, WhiteNoise et déploiement & Liens de déploiement en ligne et rapports finaux. \\
  \bottomrule
\end{tabularx}

\end{document}